\subsection{The actual complex exponential and its normalized tangent map}
\label{subsec:ftb-actual-complex-exponential}

This subsection follows the frozen source construction, before constructing the
Appell--Humbert frame.  Fix a polarised abelian variety
\(P\in\mathrm{Out.PolAV}(\operatorname{Spec}\mathbb C)\), of dimension \(g\).
We use the source's ample line bundle, affine charts and regular-function
evaluations.  The multiplication law used here is
\(\mathrm{Out.MulNSt}\): in particular, for each \(n\geq1\), multiplication
by \(n\) is a finite flat surjective morphism.  This is a previously supplied
algebraic law, not a consequence of the analytic construction below.

Write \(P(\mathbb C)\) additively.  Give it precisely the topology generated by
\[
 \{x:x\in U(\mathbb C),\ f(x)\in O\},
 \qquad U\subset P\text{ open},\quad f\in\Gamma(U,\mathcal O_P),
 \quad O\subset\mathbb C\text{ open}.                         \tag{CE1}
\]
For a map \(a:W\to P(\mathbb C)\), the source's holomorphy condition means that
\(a^{-1}U\) is open in the open set \(W\), and every
\(f\circ a\) is holomorphic on \(a^{-1}U\).  All occurrences of
holomorphic maps to algebraic varieties below have this meaning.

\begin{srlemma}[The affine and smooth charts used by the construction]
\label{lem:ftb-complex-algebraic-charts}
For an affine open \(U\subset X\) of a complex scheme, a
\(\mathbb C\)-algebra homomorphism
\(\Gamma(U,\mathcal O_X)\to\mathbb C\) determines a unique complex point of
\(U\), with these regular-function values.  If \(X\to\operatorname{Spec}\mathbb C\)
is smooth of relative dimension \(g\), then at any \(x\in X(\mathbb C)\)
there are an affine open \(U\ni x\), functions \(t_1,\ldots,t_g\), an open
neighbourhood \(W\subset\mathbb C^g\) of zero, and a map
\(\sigma:W\to X(\mathbb C)\), such that
\[
 \sigma(0)=x,\qquad t_i(\sigma(w))=w_i,                       \tag{CE2}
\]
\(\sigma\) is holomorphic, and \(\sigma(W)\) is open for (CE1).
Every derivation at \(x\) that kills the constants and all the \(t_i\)
is zero.
\end{srlemma}

\begin{proof}
The affine point is \(\operatorname{Spec}(e)\) followed by the affine-open
map.  The identities between pullback of sections and evaluation show both
its prescribed values and uniqueness; no comparison of point sets is
assumed.

Choose a standard smooth affine presentation near \(x\).  Its etale
coordinates can be translated by their values at \(x\), so that they
vanish there.  The analytic inverse used for this etale map is obtained
as follows.  Present the etale algebra over
\(\mathbb C[T_1,\ldots,T_g]\) by equally many extra variables and
relations, with invertible Jacobian in those extra variables.  If the
extra variables are \(Y\) and the relations are \(R\), apply the analytic
inverse-function theorem to
\[
             (w,y)\longmapsto (w,R(w,y)).
\]
The derivative is injective: its first component forces the variation
of \(w\) to be zero, and the invertible relation Jacobian then forces the
variation of \(y\) to be zero.  Source and target have equal finite
dimension.  Restrict the analytic inverse to \((w,0)\).
Its values satisfy every relation and hence give algebra homomorphisms
\(\sigma_w\) from the etale algebra to \(\mathbb C\).
For every algebra element, its value under \(\sigma_w\) is analytic,
because the element is represented by a polynomial in the chosen
variables.  The inverse identities give (CE2).

Here openness in (CE1) is also verified.  A finite presentation embeds
the set of algebra homomorphisms into a finite-dimensional coordinate
space.  Shrink the chart using finitely many tests
\[
       |\sigma_w(f_j)-e(f_j)|<\varepsilon .
\]
The analytic inverse contains every algebra homomorphism satisfying
these tests.  Thus its image contains the corresponding basic
neighbourhood, and after this restriction its image is such an open
chart.  A regular function on a different open set is locally a
quotient \(a/h^m\) on a basic open of \(U\).  Its value is the holomorphic
quotient \(\sigma_w(a)/\sigma_w(h)^m\), where the denominator is
nonzero.  This proves the stated holomorphy on every algebraic open.

Finally, a derivation killing the \(t_i\) kills their whole polynomial
algebra.  The etale algebra is formally unramified over that polynomial
algebra.  Apply uniqueness of lifts to the two maps into
\(\mathbb C[\epsilon]/(\epsilon^2)\), namely \(e+\epsilon D\) and \(e\).
They agree, so \(D=0\).
\end{proof}

Two further topological facts are established using these same values.
First, products of continuous, respectively holomorphic, families are
continuous, respectively holomorphic.  On affine product charts the
coordinate algebra is generated by the images of the two factor
algebras.  Addition and multiplication of their values prove the
assertion there.  Restriction to arbitrary opens is proved by
localization, using \(a/h^m\).  Applying this to the multiplication and
inverse morphisms makes \(P(\mathbb C)\) a topological group.

Second, this topological space is compact.  Here is the source's
projective-chart proof.  Choose finitely many sections of an ample
power that cover \(P\) by their nonvanishing loci, such that the ratios
of these sections generate the affine coordinate rings.  In each
chart, the homomorphisms for which all these ratios have modulus at
most one form a compact set: they are the closed zero locus of the
relations inside a finite closed unit polydisc.  At every point choose
a section of maximal modulus among the chosen sections, in any local
frame.  Its nonvanishing chart has every ratio of modulus at most one.
Consequently these finitely many compact sets cover \(P(\mathbb C)\).
The affine point correspondence in
Lemma~\ref{lem:ftb-complex-algebraic-charts} identifies their values with
the actual points throughout this argument.

\begin{srlemma}[The source's halving logarithm]
\label{lem:ftb-halving-logarithm}
Let \(V\) be a finite-dimensional complex normed vector space.
Suppose \(m\) is an analytic commutative local group law at zero, with
unit zero and associative wherever the displayed compositions are
defined near zero.  There is an analytic local change of coordinates
\(\mathcal L\), with
\[
 \mathcal L(0)=0,\quad d\mathcal L_0=1,\qquad
 \mathcal L(m(x,y))=\mathcal L(x)+\mathcal L(y)                 \tag{CE3}
\]
for sufficiently small \(x,y\).
\end{srlemma}

\begin{proof}
The unit identities and commutativity give
\[
          m(x,y)=x+y+O(\|(x,y)\|^2).
\]
The derivative of \(d(x)=m(x,x)\) is \(2\), so the analytic
inverse-function theorem gives a local inverse \(\varepsilon\), with
\[
       \varepsilon(x)=\tfrac12x+O(\|x\|^2).
\]
Associativity and commutativity give
\(d(m(x,y))=m(d(x),d(y))\).  Applying the local inverse, after shrinking
all neighbourhoods, gives
\[
       \varepsilon(m(x,y))=m(\varepsilon(x),\varepsilon(y)).   \tag{CE4}
\]
Choose \(r>0\) and \(C\geq0\) such that, for \(\|x\|<r\),
\[
 \|\varepsilon(x)\|\leq\tfrac58\|x\|,\qquad
 \|\varepsilon(x)-\tfrac12x\|\leq C\|x\|^2.
\]
All iterates remain in this ball.  Put
\(\mathcal L_j(x)=2^j\varepsilon^{\,j}(x)\).  Direct subtraction gives
the exact estimates used in the source:
\[
 \|\mathcal L_{j+1}(x)-\mathcal L_j(x)\|
       \leq2C\|x\|^2(25/32)^j,\qquad
 \|\mathcal L(x)-\mathcal L_j(x)\|
       \leq\tfrac{64}{7}C\|x\|^2(25/32)^j,                  \tag{CE5}
\]
where \(\mathcal L=\lim_j\mathcal L_j\).
The second estimate follows by summing the geometric tail.
Convergence is uniform on a smaller ball; its terms are holomorphic.
The locally uniform limit theorem therefore makes \(\mathcal L\)
holomorphic there.  Taking \(j=0\) in (CE5) gives
\(\mathcal L(x)-x=O(\|x\|^2)\), so \(d\mathcal L_0=1\).

Iterate (CE4).  If
\(\|m(x,y)-x-y\|\leq D\|(x,y)\|^2\) near zero, the difference
\[
 2^j m(\varepsilon^{\,j}(x),\varepsilon^{\,j}(y))
 -2^j\varepsilon^{\,j}(x)-2^j\varepsilon^{\,j}(y)
\]
has norm at most
\(D\|(x,y)\|^2(25/32)^j\).  It tends to zero, proving (CE3).
Another application of the inverse-function theorem gives an analytic
local inverse to \(\mathcal L\).
\end{proof}

\begin{srproposition}[The actual exponential, lattice and tangent comparison]
\label{prop:ftb-actual-complex-exponential}
There are an additive surjection
\[
                 \exp_P:\mathbb C^g\longrightarrow P(\mathbb C)
\]
and a discrete integral lattice
\(\Lambda=\ker(\exp_P)\), whose real span is \(\mathbb C^g\).
The map \(\exp_P\) is holomorphic in the regular-function sense above.
One may choose the centred coordinates \(t_i\) at the origin so that
\[
             d_0(t_1\circ\exp_P,\ldots,t_g\circ\exp_P)=1.    \tag{CE6}
\]
For each \(w\in\mathbb C^g\) there is a unique algebraic tangent vector
\(v_w:\operatorname{Spec}\mathbb C[\epsilon]/(\epsilon^2)\to P\)
at the origin such that, for every regular function defined there,
\[
          f(v_w)=f(0)+\epsilon\,d_0(f\circ\exp_P)(w).         \tag{CE7}
\]
Every algebraic tangent vector occurs uniquely in this way.
Every algebraic group endomorphism \(f:P\to P\) admits a complex
linear lift \(F\) with
\[
                         \exp_P(Fz)=f(\exp_P z).            \tag{CE8}
\]
Moreover, for every \(n\geq0\) and every regular-value-holomorphic
map \(a:\mathbb C^n\to P(\mathbb C)\) with \(a(0)=0\), there is a
map \(G:\mathbb C^n\to\mathbb C^g\), analytic at zero, with
\(G(0)=0\) and \(\exp_P(Gw)=a(w)\) for all sufficiently small \(w\).
\end{srproposition}

\begin{proof}
Use Lemma~\ref{lem:ftb-complex-algebraic-charts} at the origin.
The local operation in its coordinates is analytic by the product
calculation above.  Compose the inverse of the logarithm in
Lemma~\ref{lem:ftb-halving-logarithm} with the inverse coordinate chart.
This is a local map \(e_0\) with
\(e_0(v+w)=e_0(v)+e_0(w)\) near \((0,0)\).
It has a unique global additive extension: for \(j\) sufficiently
large set
\[
                     e(v)=2^j e_0(2^{-j}v).                \tag{CE9}
\]
Repeated local doubling shows that the right side is independent of
\(j\).  For additivity, choose one \(j\) placing \(2^{-j}v\),
\(2^{-j}w\), and \(2^{-j}(v+w)\) in the common small neighbourhood.
Uniqueness follows by applying the same rescaling to any two additive
maps that agree near zero.  Set \(\exp_P=e\).

Near zero, \(e\) is the local analytic coordinate inverse just
constructed.  Hence it is a local homeomorphism there, and translation
makes it a local homeomorphism everywhere.  Its image is an open
subgroup.  The quotient by this image is a finite group, by compactness.
Multiplication by every positive integer is surjective on \(P(\mathbb C)\):
the finite surjective algebraic multiplication morphism has a closed
point in the fibre over each complex point, and its residue field is
\(\mathbb C\).  A finite divisible abelian group is zero (multiply by
its order).  Thus \(e\) is surjective.  The local homeomorphism at zero
also makes its kernel discrete.

For completeness, the source proves full real span directly.
The image is compact, so for some \(R>0\) every point in it has a
representative \(b\) with \(\|b\|<R\): the open sets \(e(B(0,n))\)
give an increasing cover, and a finite subcover suffices.
For any \(v\) and any large positive real \(a\), choose such \(b\)
with \(e(av)=e(b)\).  Then \(av-b\in\Lambda\), and
\[
             a^{-1}(av-b)\in\operatorname{span}_{\mathbb R}\Lambda,
             \qquad \|v-a^{-1}(av-b)\|<R/a .
\]
A real linear subspace in finite dimension is closed, so the span is
all of \(\mathbb C^g\).

Holomorphy at every \(z_0\) follows from
\(e(z)=e(z_0)+e(z-z_0)\) and holomorphy near zero; the actual product
and regular-function localization calculations above apply.
The inverse logarithm has derivative the identity, giving (CE6).

For the tangent assertion, define on an affine neighbourhood \(U\)
of zero
\[
 J_w(a)=a(0)+\epsilon\,d_0(a\circ e)(w).
\]
The product rule makes \(J_w\) a \(\mathbb C\)-algebra homomorphism;
the affine point correspondence over the dual numbers gives \(v_w\).
For an arbitrary tangent \(v\), take
\(w_i\) to be the \(\epsilon\)-coefficient of \(t_i(v)\).
The difference of the two corresponding derivations kills all
constants and the \(t_i\), so it is zero by the last assertion of
Lemma~\ref{lem:ftb-complex-algebraic-charts}.  This proves existence
and uniqueness on \(U\).  Localization to basic opens and restriction
of the evaluation homomorphisms proves (CE7) for every regular
function defined at the origin.

For the general local lifting assertion, write \(\varphi\) for the
origin coordinate chart used above.  The map \(\varphi\circ a\) is
analytic at zero by the regular-value definition and has value zero.
Put \(G=\mathcal L\circ\varphi\circ a\) near zero, and extend it
arbitrarily outside this neighbourhood.  The local inverse identities
give \(G(0)=0\) and \(e\circ G=a\) near zero.
Apply this assertion to \(a=f\circ e\) to obtain a holomorphic
\(G:\mathbb C^g\to\mathbb C^g\) near zero, with \(G(0)=0\) and
\(e(Gz)=f(ez)\).
The difference \(G(2z)-2G(z)\) belongs to the discrete kernel and
tends to zero.  It is therefore zero for all sufficiently small \(z\).
For a fixed \(z\), the sequence \(2^jG(2^{-j}z)\) is eventually
constant, while differentiability makes its limit \(dG_0(z)\).
Applying \(e\) and repeatedly doubling the identity
\(e(G(2^{-j}z))=f(e(2^{-j}z))\) proves (CE8) with \(F=dG_0\).
\end{proof}

\paragraph{Exact source scope and dependency ledger.}
The output here is the exponential part of
\srcid{AbelianExp.exists\_holUnif}, including its extra tangent-coordinate
and local lifting assertions; it initially uses the trivial line bundle
with constant frame \(1\).  It does not yet give the frame of the
polarising bundle, its positive Hermitian form, the Weil pairing formula,
or section--theta correspondence.  Those remain the subsequent,
separate steps of \srcid{thetaUnifAV\_of\_laws}.

All paths in this ledger are under
\srcid{src/Development/AbelianVarietyTheory/Proofs/} at tag
\srcid{alpoge-source-7cbf7a0c}.
The affine point proof is
\srcid{ToricGeometry/ExistsUniquePointOfRingHom.lean}
(blob \srcid{b4794489f6856874c710b89c8f1d3f65f8437af9}).
The analytic etale inverse is
\srcid{HomogeneousIdeals/ExistsOpenPartialHomeomorphAnalytic.lean}
(\srcid{3574b51a9512022e9702954181112212d94db05f});
its finite-test form is
\srcid{ToricGeometry/ExistsTestFunctionsSubsetOfMemNhds.lean}
(\srcid{5ba21a914397f864d2e7b2660a78dd12a9d8e243}), and the smooth
algebraic charts and derivation uniqueness are
\srcid{ToricGeometry/ExistsEtaleCoordsVanishing.lean}
(\srcid{1c87b47bf131dbe99f0958c2abc699184d7a633a}).
The topological group and compactness proofs are respectively
\srcid{ToricGeometry/IsTopologicalGroupComplexPoints.lean}
(\srcid{7e9ca75c5be097e97403b219612289baa50acd6e}) and
\srcid{ToricGeometry/CompactSpacePointsPolAV.lean}
(\srcid{b90f408674114e76130f561572900076bdffdcf6}).
Product holomorphy is
\srcid{ToricGeometry/HolFamilyOfIsPullback.lean}
(\srcid{c343c07f96ec7ec7d2e61d8735cbf9b1b16cf273}).

The quantitative halving limit and extension (CE9) are
\srcid{HomogeneousIdeals/TendstoUniformlyOnDoublingLogOfContraction.lean}
(\srcid{69e6b8b788ca80ded9032668f76eec74e50ed461}).
The local group-law identities, analytic inverse logarithm, discrete
kernel, full span and compact/divisible surjectivity are
\srcid{HomogeneousIdeals/IsLocalHomeomorphAddMonoidHomOfEventuallyEq.lean}
(\srcid{085296514fbf6aaab8595e9593f38e6f4ecc8efc}).
Their actual application to \(P\) is, despite its file name,
\srcid{ToricGeometry/ContinuousOnValue.lean}
(\srcid{967923ce3ac55f13e79463c63e672671f9edcf37}),
the declaration \srcid{AbelianExp.exists\_holUnif}.
The tangent construction is
\srcid{ToricGeometry/TangentBijectionOfValues.lean}
(\srcid{ae10755ee2ab73933551218a57128e702c5cd2c4});
the endomorphism argument is
\srcid{ToricGeometry/ExistsLinearLift.lean}
(\srcid{fb716479e1ebeb2427d806db980225401efb1c21}).
The fixed ordering of these calls is
\srcid{ToricGeometry/ThetaUnifAVOfLaws.lean}
(\srcid{3e496ea337ef9e519bbf43dfe91ef52009f6a960}).
The standard analytic inverse-function and locally uniform limit
theorems are used with the finite-dimensional complex and open-domain
hypotheses displayed above.
